\section{Procedura Git}

\subsection{Setup iniziale - Clonare la repository}

Per iniziare a lavorare al progetto è necessario clonare la repository. Eseguire dal terminale:

\begin{lstlisting}
git clone https://github.com/Idrorepellente/quantsys_institutional.git
cd quantsys_institutional
\end{lstlisting}

Successivamente, creare un virtual environment e installare le dipendenze:

\begin{lstlisting}
python -m venv venv
# Su Windows PowerShell:
.\venv\Scripts\Activate.ps1
# Su Linux/macOS:
source venv/bin/activate

pip install -r requirements.txt
\end{lstlisting}

Per i dettagli completi sul setup dell'ambiente (configurazione \texttt{.env}, scaricamento delle cache dati) si rimanda al Manuale Utente.

\subsection{Creazione di una nuova feature}

Il flusso operativo per la creazione di una nuova feature si articola nei seguenti step.

\subsubsection{Step 1: Identificare o creare la issue}

Ogni nuova feature deve essere collegata a una issue su GitHub. Se non esiste già, aprire una nuova issue specificando:

\begin{itemize}
  \item \textbf{Titolo}: descrittivo e conciso;
  \item \textbf{Descrizione}: motivazione della feature, criteri di accettazione, eventuali riferimenti a letteratura o discussioni precedenti;
  \item \textbf{Labels}: assegnare label appropriati (es. \texttt{enhancement}, \texttt{bug}, \texttt{strategy}, \texttt{docs});
  \item \textbf{Assignee}: indicare il responsabile dello sviluppo.
\end{itemize}

Annotare il numero della issue (es. \#42), che verrà utilizzato per il nome del branch.

\subsubsection{Step 2: Aggiornare il branch \texttt{main} locale}

Prima di creare il nuovo branch, assicurarsi di avere l'ultima versione di \texttt{main}:

\begin{lstlisting}
git checkout main
git pull origin main
\end{lstlisting}

\subsubsection{Step 3: Creare il branch dedicato}

Creare il branch utilizzando la convenzione di naming descritta nel Capitolo \ref{chap:processo_supporto}:

\begin{lstlisting}
git checkout -b feature/42-nome-descrittivo
\end{lstlisting}

\subsubsection{Step 4: Sviluppo locale}

Effettuare le modifiche necessarie. Durante lo sviluppo è opportuno effettuare commit incrementali, ciascuno con un messaggio chiaro:

\begin{lstlisting}
git add .
git commit -m "Aggiungi calcolo ATR(14) all'indicator pool"
\end{lstlisting}

In fase di sviluppo, se si tocca codice rilevante per la pipeline di backtest, è raccomandato:

\begin{itemize}
  \item Eseguire i test unitari della cartella \texttt{tests/};
  \item Effettuare almeno un walk-forward completo della strategia o componente coinvolto, verificando che i numeri OOS non degradino rispetto allo stato precedente.
\end{itemize}

\subsection{Caricamento dei cambiamenti}

\subsubsection{Step 1: Verificare lo stato dei file}

Prima di eseguire il push, verificare lo stato del working directory:

\begin{lstlisting}
git status
\end{lstlisting}

\subsubsection{Step 2: Aggiungere i file modificati}

\begin{lstlisting}
git add .
\end{lstlisting}

Se si desidera aggiungere solo specifici file:

\begin{lstlisting}
git add core/exit_policy.py strategies/momentum/strategy.py
\end{lstlisting}

\subsubsection{Step 3: Eseguire il commit}

\begin{lstlisting}
git commit -m "Aggiungi exit_policy_type al StrategyDescriptor"
\end{lstlisting}

\subsubsection{Step 4: Eseguire il push}

\begin{lstlisting}
git push origin feature/42-nome-descrittivo
\end{lstlisting}

Al primo push di un branch nuovo è necessario specificare l'upstream:

\begin{lstlisting}
git push --set-upstream origin feature/42-nome-descrittivo
\end{lstlisting}

\subsection{Sincronizzazione con \texttt{main}}

Se il branch di lavoro è "in vita" da molti giorni, è opportuno periodicamente sincronizzarlo con \texttt{main} per evitare conflitti al merge:

\begin{lstlisting}
git checkout main
git pull origin main
git checkout feature/42-nome-descrittivo
git merge main
\end{lstlisting}

In caso di conflitti:

\begin{enumerate}
  \item Aprire i file con conflitti nell'editor (sono evidenziati da \texttt{git status});
  \item Risolvere manualmente i conflitti scegliendo cosa mantenere;
  \item Eseguire \texttt{git add <file>} sui file risolti;
  \item Concludere il merge con \texttt{git commit};
  \item Effettuare il push.
\end{enumerate}
